% New English translation from checked cote 25, site 104 (PDF page 105).
% AI-translated by GPT-6 Astra (OpenAI, ChatGPT Pro mode).
\scanpage{25/104}
\mspage{7}

\uline{ms} \uline{constant} rel. to S. 2) n=1, then we must take the homomorphisms
\var{\xout{everywhere surjec}}{}{typewriter strike-out}{ega25_25_104-v0001} $\underline{E}'\hand{\to} \underline{E}\otimes\underline{M}$ that are \uline{surjective}, as one verifies
immediately, and the corresponding homomorphism \unsure{$g$|$f$}: $P\hand{\to} P'$ is precisely the
composite $P(\underline{E})\hand{\xrightarrow{\sim}}P(\underline{E}\otimes\underline{M})\hand{\to} P(\underline{E}')$, where the first hom is the canonical
isomorphism considered in Chap II, and the second the canonical closed immersion
obtained from the epimorphism $\underline{E}'\hand{\to}\underline{E}\otimes\underline{M}$. If we call \uline{linear}
the homs from P to P' described as such a composite, we therefore see
that the morphisms $g:P\hand{\to} P'$ of degree 1 are exactly those that
are \uline{linear}.

To finish, let us determine the isomorphisms of P with P':
\uline{Theorem} 1.6. Let S be a prescheme, $P=P(\underline{E})$ and $P'=P(\underline{E}')$ two projective
bundles over S\var{\inserted{defined \unsure{by|by} $\underline{E}$, $\underline{E}'$ loc. free of finite type.}}{ defined \unsure{by|by} $\underline{E}$, $\underline{E}'$ loc. free of finite type.}{insertion}{ega25_25_104-v0002}. Then every S-isomorphism $\unsure{g|f}$: $P\hand{\xrightarrow{\sim}}P'$ can be defined
as the composite $P(\underline{E})\hand{\xrightarrow{\sim}}P(\underline{E}\otimes\underline{M})\hand{\to} \underline{P}(\underline{E}')$, where\var{\inserted[typed]{$\underline{M}$ is an invertible Module on S,}}{ $\underline{M}$ is an invertible Module on S,}{insertion}{ega25_25_104-v0003} the first hom is the can
isom of Chap II, and the \unsure{second|second} is the \var{\xout{canonical}}{}{typewriter strike-out}{ega25_25_104-v0004} isomorphism obtained
from an isomorphism u: $E'\hand{\xrightarrow{\sim}}\underline{E}\otimes\underline{M}$. When the fibers of \unsure{P|p} are
nonempty (resp. of dim $\hand{\geq} 1$) then $\underline{M}$ (resp. the pair $(\underline{M},u)$) is determined
up to unique isomorphism in terms of g.

By the preceding considerations, we reduce to proving
that g is of degree 1, which reduces us to the case where S is the spectrum of a
field,\var{\inserted[typed]{and of course we can suppose dim $P\hand{\geq} 1$.}}{ and of course we can suppose dim $P\hand{\geq} 1$.}{insertion}{ega25_25_104-v0005} But note that $\underline{O}_P(1)$ is then characterized intrinsically (i.e.
independently of how P has been realized as a projective bundle)
\unsure{as|as} the generator of Pic(P) (among the two generators $\underline{O}_P(1)$ and
$\underline{O}_{\unsure{P|E}}(-1)$) that is \uline{ample}; consequently if g: $P\hand{\to} P'$ is an isomorphism,
then $g^{*}(\underline{O}_{P'}(1))$ is isomorphic to $\underline{O}_P(1)$, and we win. -- In a local
% NOTE: the right edge of the source scan cuts off the end of "local[e]".
and less learned form, we can state:
\uline{Corollary} 1.7. Under the conditions of 1.6., every\var{\inserted{S-}}{ S-}{insertion}{ega25_25_104-v0007}isomorphism $g:P\hand{\to} P'$
can be described, in a neighborhood of every $s\hand{\in} S$, using an isomorphism
u: $\underline{E}\hand{|}U \hand{\simeq} E'\hand{|}U$, the latter being well determined mod mult by an element
\tnote{25-104-1}{The live prose alternatives pr or par, deuxième or dauxième, and the final e or a in comme have coincident English renderings; their uncertainty marks remain. The canceled uncertain \unsure{e|x} after comme has no English spelling effect. The source abbreviations can isom and mod mult mean canonical isomorphism and modulo multiplication. The underlined P in the theorem's composite and the plain E-prime in its isomorphism remain as written.}
\var{\margin{104}}{\placedmargin{104}}{marginal note placed}{ega25_25_104-v0008}
